\section{Discussion}
\label{sec:discussion}

% Paragraph 1 -- Clinical implications
The operating-point metrics reported in Table~\ref{tab:confusion} have direct implications for population-level screening.
A sensitivity of 85.2\% means that the model identifies the large majority of individuals whose fasting plasma glucose (FPG) exceeds the diabetic threshold, while the 98.0\% negative predictive value (NPV) indicates that patients receiving a negative screen can be reassured with high confidence.
The accompanying false-positive rate of 37.5\% is a deliberate consequence of the recall-oriented threshold: in a screening context, a false positive triggers only a confirmatory blood glucose test, not treatment, making the downstream cost low compared with a missed diagnosis.
Traditional questionnaire-based risk scores such as ADA, IDRS, and FINDRISC have demonstrated lower sensitivity than laboratory-informed models when applied to population screening~\cite{Doddamani2021comparative}, and the USPSTF recommends screening adults aged 35--70 with overweight or obesity~\cite{USPSTF2021screening}, a criterion that misses metabolically unhealthy normal-weight individuals captured by our model's lab-based features.
As Deng et~al.\ emphasize, translating a high area under the receiver operating characteristic curve (AUC) into clinical benefit requires attention to calibration, threshold choice, and deployment context~\cite{Deng2025highauc}; our Platt-calibrated probabilities and recall-targeted threshold address these concerns explicitly.

% Paragraph 2 -- Tree models vs deep learning on tabular data
CatBoost achieved an AUC of 0.817, matching the PyTorch Tabular network (0.816) and the sklearn multilayer perceptron (0.812) within 0.005 AUC (Table~\ref{tab:model-comparison}).
This outcome is consistent with the NeurIPS finding of Grinsztajn et~al.\ that tree-based models match or outperform deep learning on typical tabular datasets~\cite{Grinsztajn2022trees}.
Beyond raw discrimination, CatBoost offers practical advantages for clinical deployment: it provides exact Shapley values through its native TreeSHAP implementation rather than the approximate values produced by model-agnostic methods, it handles missing values internally without requiring a separate imputation step, and it trains and infers substantially faster than neural alternatives.
For a screening tool that must be transparent to clinicians, the alignment of SHAP feature rankings with established risk factors---age, liver enzyme ratios, waist-to-height ratio, lipid ratios (Table~\ref{tab:shap})---is not merely convenient but necessary for building trust and facilitating adoption.

% Paragraph 3 -- Missingness as signal
Approximately 66\% of records in the NHIS dataset lack a lipid panel, reflecting the optional nature of this test in Korean biennial health examinations.
CatBoost handles these missing values natively through learned default split directions, eliminating the need for imputation.
We additionally engineered nine explicit binary missingness flags, yet all exhibited effectively zero SHAP importance (maximum $1.7 \times 10^{-7}$), indicating that the algorithm's internal mechanism fully subsumes the information these indicators were intended to capture.
This finding is consistent with prior work showing that explicit missingness indicators provide limited additional value when the base learner already models absence~\cite{Singh2021missingness, Sisk2023imputation}.
The practical consequence is favorable: the model degrades gracefully in clinical settings that do not routinely order lipid panels, and no external imputation pipeline is required at inference time.

% Paragraph 4 -- Three-tier design and practical deployment
The feature-set ablation in Table~\ref{tab:ablation} reveals an asymmetric cost--benefit structure.
Adding non-lipid laboratory values to the core demographic and anthropometric features produces the largest AUC gain (+0.033, from 0.775 to 0.808), whereas appending the full lipid panel yields only a marginal further improvement (+0.009, to 0.817).
This result carries policy implications: effective T2D screening is achievable without cholesterol tests, reducing per-patient cost and broadening coverage to settings where lipid panels are unavailable.
Even the core-only model (15~features, AUC 0.775) provides useful risk stratification in resource-constrained environments.
The alignment of top SHAP features---age, liver enzymes, waist circumference, lipid ratios---with well-established clinical risk factors supports the model's face validity and suggests that the learned patterns reflect genuine pathophysiological relationships rather than dataset artifacts.
We operationalized these findings in a deployed web-based screening tool that maps calibrated probabilities to a four-tier risk stratification (low, moderate, high, very high), providing clinicians with an actionable output rather than a raw score.
However, the model was trained exclusively on South Korean NHIS data, and its generalizability to other populations---with different dietary patterns, genetic backgrounds, and healthcare delivery structures---remains to be established through external validation.
